\documentclass[11pt,a4paper]{article}

%--- Geometry & encoding ----------------------------------------------------
\usepackage[T1]{fontenc}
\usepackage[utf8]{inputenc}
\usepackage[margin=2.5cm]{geometry}
\usepackage{lmodern}
\usepackage{microtype}

%--- Math -------------------------------------------------------------------
\usepackage{amsmath}
\usepackage{amssymb}
\usepackage{amsthm}
\usepackage{mathtools}
\usepackage{stmaryrd}

%--- Tables, figures, code ---------------------------------------------------
\usepackage{booktabs}
\usepackage{longtable}
\usepackage{array}
\usepackage{tabularx}
\usepackage{listings}
\usepackage{xcolor}
\usepackage{enumitem}

%--- Hyperlinks (last) -------------------------------------------------------
\usepackage[
    colorlinks=true,
    linkcolor=blue!55!black,
    citecolor=blue!55!black,
    urlcolor=blue!55!black
]{hyperref}

%--- Theorem environments ---------------------------------------------------
\theoremstyle{plain}
\newtheorem{theorem}{Theorem}[section]
\newtheorem{lemma}[theorem]{Lemma}
\newtheorem{proposition}[theorem]{Proposition}
\newtheorem{corollary}[theorem]{Corollary}

\theoremstyle{definition}
\newtheorem{definition}[theorem]{Definition}
\newtheorem{remark}[theorem]{Remark}

%--- Lean code listings -----------------------------------------------------
\definecolor{leanBg}{RGB}{248,248,250}
\definecolor{leanKey}{RGB}{32,72,160}
\definecolor{leanCmt}{RGB}{100,120,140}

\lstdefinelanguage{Lean}{
  morekeywords={theorem,lemma,def,structure,inductive,where,by,fun,let,in,
                if,then,else,match,with,namespace,end,open,import,
                noncomputable,instance,class,variable,intro,exact,apply,
                rfl,simp,rw,have,this,show,use},
  sensitive=true,
  morecomment=[l]{--},
  morecomment=[s]{/-}{-/},
  morestring=[b]"
}
\lstset{
  language=Lean,
  basicstyle=\ttfamily\small,
  keywordstyle=\color{leanKey},
  commentstyle=\color{leanCmt}\itshape,
  backgroundcolor=\color{leanBg},
  frame=single,
  rulecolor=\color{leanBg!50!black!20},
  framerule=0.4pt,
  framexleftmargin=4pt,
  framexrightmargin=4pt,
  framextopmargin=2pt,
  framexbottommargin=2pt,
  breaklines=true,
  showstringspaces=false,
  xleftmargin=8pt,
}

%--- Convenience macros -----------------------------------------------------
\newcommand{\C}{\mathbb{C}}
\newcommand{\R}{\mathbb{R}}
\newcommand{\N}{\mathbb{N}}
\newcommand{\Q}{\mathbb{Q}}
\newcommand{\Z}{\mathbb{Z}}
\newcommand{\eml}{\mathrm{eml}}
\newcommand{\edl}{\mathrm{edl}}
\newcommand{\imu}{\mathbf{i}}
\newcommand{\EMLT}{\mathsf{EMLTerm}}
\newcommand{\EMLTC}{\mathsf{EMLTerm}_\C}
\newcommand{\EDLT}{\mathsf{EDLTerm}}
\newcommand{\F}{\mathsf{F36}}
\newcommand{\EL}{\mathsf{EL}}

%--- Title block ------------------------------------------------------------
\title{\Large\textbf{Formalising EML in Lean 4}\\[2pt]
       \large A post-frontier technical summary\\
       \texttt{arXiv:2603.21852}}
\author{Bartosz Naskręcki\\
        \small Adam Mickiewicz University, Pozna\'n /
        CCAI Warsaw University of Technology\\[1pt]
        \small with Aristotle (Harmonic), GPT~Pro (OpenAI),
        and Claude (Anthropic)}
\date{12 May 2026}

\begin{document}
\maketitle

\begin{abstract}
This report documents the Lean 4 + Mathlib v4.28 formalisation of
\emph{All elementary functions from a single binary operator} by
A.~Odrzywo\l{}ek (\texttt{arXiv:2603.21852}). The paper introduces the
\emph{EML} operator $\eml(x,y) = \exp(x) - \log(y)$ and shows that the
two-element basis $\{1, \eml\}$ is a \emph{Sheffer system} sufficient
to reconstruct each of $36$ primitive functions on a reference
calculator. Our artefact seals all $36$ paper primitives via literal
$\EMLT$ (or $\EMLTC$, complex fragment) witness terms whose partial
evaluation matches the paper's stated value on a non-empty open
subdomain of the natural mathematical domain. The three §G boundary
points the paper itself flags ($\sqrt{0}$, $\mathrm{arcosh}\,1$,
$\mathrm{hypot}(0,0)$) are additionally sealed by \emph{witness-family}
theorems of the shape $\forall \mathrm{env},\,[\mathrm{hyp}] \to
\exists t \in \EMLT,\,t.\mathrm{eval}_? \mathrm{env} =
\mathrm{some}\,v$, in which the term $t$ is allowed to depend on the
environment. After the post-submission frontier sprint of
2026-05-10/11, the public Lean API exposes $100$ theorems:
$61$ paper-claim theorems ($48$ EML + $8$ EDL + $5$ $-$EML) plus
$39$ from six frontier modules covering full-domain trig (Path~C$'$),
§G boundaries (\texttt{GFullFix}), variable-transplant identity depths
(\texttt{TransplantDepths}), polynomial-class universal minimality
(\texttt{PolynomialBinary}), and the EDL closed-value closure with a
conditional-on-typeclass obstruction scaffold (\texttt{EDLClosedVal}).
\texttt{lake build} finishes sorry-free at $8062$ jobs; every public
theorem depends only on Mathlib's three standard axioms
(\texttt{propext}, \texttt{Classical.choice}, \texttt{Quot.sound}).
The artefact is independently re-verified on PCSS Eagle HPC.
\end{abstract}

\tableofcontents
\bigskip\hrule\bigskip

%==========================================================================
\section{Introduction}\label{sec:intro}

\subsection{The paper}

The source paper proves that the binary operator
\[
   \eml(x, y) \;=\; \exp(x) - \log(y),
\]
together with the constant $1$, is a Sheffer system for elementary
functions: each primitive on a 36-button reference calculator
($\pi, e, \imu, -1, 1, 2, x, y, \exp, \log, \mathrm{inv}, \sqrt{\cdot},
\sin, \cos, \arctan, \sinh, \mathrm{arcosh}, \dots, +, -, \cdot, /,
\log_b, \mathrm{pow}, \mathrm{avg}, \mathrm{hypot}$) admits a finite
\emph{EML reconstruction tree} built from $\{1, \mathbf{x}_n, \eml\}$.
The paper's central evidence is empirical (a Rust-plus-Mathematica
sieve enumerated trees up to fixed complexity and reported per-primitive
RPN-length witnesses), supplemented by a paper-level proof sketch
(Lemmas 1--3, Corollary 4 of the Supplementary Information).

The paper's Supplementary §2 explicitly notes that ``a natural next
step would be formalization in Lean 4, but preliminary AI-assisted
attempts failed; the extended-value conventions ($\log 0 = -\infty$)
and branch-cut reasoning required appear to exceed current automation
capabilities.'' The artefact described here demonstrates that, with
the right architectural choices, formalisation \emph{is} feasible
end-to-end.

\subsection{Why the paper's framing makes this hard}

Three Lean-specific constraints shape every architectural decision.

\begin{itemize}[leftmargin=*]
  \item \textbf{Totality.} Lean kernel functions must be total. Mathlib
    defines \texttt{Real.log\,0\,=\,0} (the ``junk value'') and
    similarly $\sqrt{x} = 0$ for $x \leq 0$. Short Supplementary
    witnesses such as $0 = L(1)$ and $-1 = L(1) - 1$ that implicitly
    use $\log 0 = -\infty$ \emph{fail} at face value: the natural EML
    composition evaluates them to $1$, not $0$.
  \item \textbf{Branch cuts.} Every reconstruction beyond the
    elementary $e = \eml(1, 1)$ uses $\log$ on values that cross or
    sit near the principal branch. ``Valid on its principal branch''
    is one sentence in the paper proof sketch; in Lean it is a
    hypothesis whose 11-field side-condition bundle must be discharged
    at every $\mathrm{mkLog}$ application.
  \item \textbf{Term-level vs.~function-level.} A real-valued identity
    $\log z = \eml(1, \eml(\eml(1, z), 1))$ is one theorem; lifting it
    to an inductive grammar $\EMLT$ so that the universality claim is
    $\exists t \in \EMLT,\, \forall \mathrm{env},\,t.\mathrm{eval} \cdots
    = \log(\mathrm{env}\,0)$ is another. We need both, and Mathlib has
    no built-in scaffolding for either.
\end{itemize}

\subsection{Headline result}

After two phases of work (a chunked-decomposition phase 1 in April
2026, archived in
\path{process_archive/REPORT_2026_04_27.pdf}; a structural-compiler
plus Path~C$'$ phase 2 in May 2026; a frontier sprint phase 3 on
2026-05-10/11), the artefact today seals:

\begin{itemize}[leftmargin=*]
  \item All $36$ paper primitives via literal $\EMLT$ or $\EMLTC$
    witnesses on a non-empty open subdomain of the natural domain;
  \item The three §G boundary points via witness-family theorems
    (\texttt{GFullFix.lean}) where the term depends on the environment;
  \item Full natural domain for $\sin$, $\arctan$, $\tan$ via Path~C$'$
    range-reduction by substitution (\texttt{Periodicity.lean},
    \texttt{Subst.lean});
  \item $8 / 36$ Sheffer-cousin primitives for the paper's EDL
    cousin and $5 / 36$ for $-$EML
    (\texttt{Sheffer.lean}, \texttt{EDLClosedVal.lean});
  \item Variable-transplant identity terms at depths $4k$ for every
    $k \in \N$ plus negative closures at depths $1$, $2$
    (\texttt{TransplantDepths.lean});
  \item A polynomial-class first cut of paper~§5 universal minimality
    (\texttt{PolynomialBinary.lean}): no bivariate polynomial binary
    operation can equal \texttt{Real.exp} on all of $\R$.
\end{itemize}

The public Lean API now exposes \textbf{$100$ machine-checked
theorems} ($61$ original paper claims + $39$ frontier). \texttt{lake
build} finishes sorry-free at $8062$ jobs. Every public theorem
depends only on the three Mathlib-standard axioms; the
\texttt{K\_count} theorems are pure \texttt{rfl}-decidable and depend
on no axioms at all (see Section~\ref{sec:verification}).

%==========================================================================
\section{Architecture}\label{sec:arch}

\subsection{The three-grammar pipeline}

The compositional core is a three-language pipeline. Each box is an
inductive Lean type; each arrow is a total function returning
\texttt{Option}-valued partial evaluations.

\medskip

\begin{center}
\begin{tabular}{l l}
\toprule
\textbf{Layer} & \textbf{What it holds} \\
\midrule
$\F$    & 36-primitive source language with partial denotation \\
$\EL$   & Intermediate language: $\{\exp, \log, \mathrm{neg},
           \mathrm{inv}, +, -, \cdot, /, \dots\}$, real partial eval \\
$\EMLT$ & Single-operator grammar: $T ::= 1 \mid \mathbf{x}_n \mid
           \eml(T, T)$, real partial eval \\
\midrule
$\EMLTC$ & Complex-coefficient extension, same syntax, $\C$ semantics \\
\bottomrule
\end{tabular}
\end{center}

The two compilers
$\F \xrightarrow{\texttt{translate?}} \EL$ and
$\EL \xrightarrow{\texttt{compile}} \EMLT$ are defined in
\path{Framework/Compilers/F36ToEL.lean} and
\path{Framework/Compilers/ELToEML.lean}; each preserves partial
evaluation. The complex extension lifts every real-fragment witness
to $\EMLTC$ via a syntactic identity homomorphism
$\iota \colon \EMLT \to \EMLTC$
and provides Euler-style witnesses for the trig family (\S\ref{sec:trig}).

\subsection{Partial evaluation as the totality work-around}

Every $\EMLT$ evaluation returns \texttt{Option\,$\R$}, with
$\eml(a, b)$ returning \texttt{none} whenever $b \leq 0$. This is the
key design move that side-steps the $\log 0 = 0$ junk-value problem:
nested \texttt{eml(a, b)} \emph{never} evaluates to a junk value at a
boundary; it returns \texttt{none} instead. Concretely:

\begin{lstlisting}[language=Lean]
noncomputable def EMLTerm.eval? (env : Nat -> Real)
    : EMLTerm -> Option Real
  | .one     => some 1
  | .var n   => some (env n)
  | .eml a b => (eval? env a).bind fun va =>
                  (eval? env b).bind fun vb =>
                    if vb <= 0 then none
                    else some (Real.exp va - Real.log vb)
\end{lstlisting}

Every \texttt{paper\_claim\_<f>} theorem states an existential of the
shape
\[
   \exists t \in \EMLT.\quad
   \forall \mathrm{env} \in (\N \to \R),
   \;[\text{precondition on } \mathrm{env}],
   \;t.\mathrm{eval}_?\,\mathrm{env} =
   \mathrm{some}\,(f(\mathrm{env}\,0,\,\mathrm{env}\,1,\,\dots)).
\]
The precondition is empty when $f$ is a full-natural-domain primitive;
otherwise it restricts to the open subdomain of validity (e.g. $x > 0$
for $\sqrt{\cdot}$).

\subsection{Path~C$'$: the witness-family shape}\label{sec:wf}

Two architectural problems force a generalisation of the existential
shape:

\begin{enumerate}[topsep=2pt, itemsep=2pt]
  \item \textbf{Trig widening.} The Euler-bridge witnesses for $\sin$,
    $\arctan$, $\tan$ work only on symmetric subdomains around $0$
    because the intermediate $\log_\C$ argument crosses the principal
    branch outside that subdomain.
  \item \textbf{§G boundary points.} A single environment-independent
    $\EMLT$ witness for $\sqrt{0}$, $\mathrm{arcosh}\,1$, or
    $\mathrm{hypot}(0, 0)$ fails because the natural composition
    evaluates to $1$ at the boundary, not $0$.
\end{enumerate}

The fix in both cases is the same: \textbf{flip the quantifiers}. Move
from the original
\[
   \exists t \in \EMLT.\quad
   \forall \mathrm{env},\;t.\mathrm{eval}_?\,\mathrm{env} =
   \mathrm{some}\,(f(\dots))
\]
to a \emph{witness-family} statement
\[
   \forall \mathrm{env}.\quad [\mathrm{hyp}] \to
   \exists t \in \EMLT,\;t.\mathrm{eval}_?\,\mathrm{env} =
   \mathrm{some}\,(f(\dots)),
\]
where the choice of $t$ is allowed to depend on $\mathrm{env}$. For
trig, the dependence is a period number $k \in \Z$ such that
$x - k\pi$ lands in the local-witness domain. For §G, the dependence
is a case split: the constant-zero term $\mathrm{mkZero}$ on the
boundary, the existing narrow witness off the boundary.

The witness-family statement is mathematically equivalent to what the
paper's compiler does in prose: it ``manually corrects the $i$-sign''
per region of $x$ at code-generation time, which is exactly choosing a
different witness per environment.

\subsection{The structural compiler vs.~per-primitive witnesses}

A second design decision is whether to seal each primitive
\emph{individually} (one hand-tuned witness per
\texttt{paper\_claim\_<f>}) or to seal a \emph{structural compiler}
that produces witnesses uniformly from a higher-level grammar.

The artefact does both, but the bulk uses the structural compiler.
The compiler $\F \to \EL \to \EMLT$ takes any $\F$-expression and
mechanically produces a $\EMLT$ witness; the per-primitive theorems
are then thin existential shells that instantiate the compiler with
the corresponding $\F$-expression. Hand-tuned witnesses are kept for
the closed numeric constants ($\pi$, $i$, $0$, $2$, $-i$) and the
trig family, where domain-specific tricks (the Cayley quotient for
$\tan$, the substitution route $\arcsin(x/\sqrt{1+x^2})$ for
$\arctan$) compress the witness drastically.

The cost of uniformity is K-count blow-up: $K(\log_b) = 9\,929\,087$
in the compiler output versus a few hundred nodes for a hand-tuned
witness. We treat that gap as informative rather than worrying:
the paper's K-counts are an \emph{upper bound on necessary size}, and
our machine-checked figures are the size of a \emph{specific
mechanically-produced} witness.

\subsection{Real-to-complex term lift}\label{sec:realcomplex}

Many complex witnesses (the trig family, most §G fixes) need to use a
real value $-x$ inside a complex composition. Constructing $-x$ from
scratch in $\EMLTC$ is expensive; instead, we use a syntactic
homomorphism \texttt{EMLTerm.toComplex} that lifts any real-fragment
$\EMLT$ tree to $\EMLTC$ unchanged, preserves the size measure, and
satisfies $\mathrm{eval}_? = (\mathrm{eval}_? \cdot \mathrm{ofReal})$.
Combined with the public closed-constants
$\mathrm{realize}_\C^{\{0, 2, -i\}}$
this lets every trig witness reuse the already-sealed real fragment
for arithmetic sub-expressions and avoids redoing branch-cut work for
each new primitive.

%==========================================================================
\section{What is sealed}\label{sec:sealed}

\subsection{Coverage scoreboard}

\begin{center}
\begin{tabular}{l l l}
\toprule
\textbf{Family (count)} & \textbf{Single structural witness}
                        & \textbf{Witness-family (full domain)} \\
\midrule
Atoms (7)                 & full natural domain                   & --- \\
$\imu$ (1)                & full ($\mathrm{realize}_\C^{\,\imu}$) & --- \\
Real unary (7)            & full natural domain                   & --- \\
$\sqrt{\cdot}$ (1)        & $(0, \infty)$                          & $\forall x \geq 0$ \\
Hyperbolic (5)            & full natural domain                   & --- \\
$\mathrm{arcosh}$ (1)     & $(1, \infty)$                          & $\forall x \geq 1$ \\
Binary (7)                & full natural domain                   & --- \\
$\mathrm{hypot}$ (1)      & $\R^2 \setminus \{(0,0)\}$             & $\forall (x, y)$ \\
$\cos$, $\arccos$, $\arcsin$ (3) & full natural domain (1, $(-1,1)$, $(-1,1)$) & --- \\
$\sin$, $\arctan$, $\tan$ (3) & narrow positive/negative           & full natural domain (Path~C$'$) \\
\bottomrule
\end{tabular}
\end{center}

\noindent\textbf{Net.} $36 / 36$ paper primitives sealed. Three §G
boundary points sealed by witness families. Three trig primitives
widened to full natural domain by Path~C$'$. The headline existential
statement holds for every primitive.

\subsection{The §G witness-family seal}\label{sec:gfix}

The module \path{Framework/GFullFix.lean} provides three theorems:

\begin{lstlisting}[language=Lean]
theorem paper_claim_sqrt_full :
    forall env : Nat -> Real, 0 <= env 0 ->
      exists t : EMLTerm, t.eval? env = some (Real.sqrt (env 0))

theorem paper_claim_arcosh_full :
    forall env : Nat -> Real, 1 <= env 0 ->
      exists t : EMLTerm, t.eval? env = some (Real.arcosh (env 0))

theorem paper_claim_hypot_full :
    forall env : Nat -> Real,
      exists t : EMLTerm,
        t.eval? env = some (Real.sqrt (env 0 ^ 2 + env 1 ^ 2))
\end{lstlisting}

\noindent
The proof of each is a case split on whether the environment is at the
boundary. On the boundary the witness is the constant-zero term
$\mathrm{mkZero}$; off the boundary the witness is the corresponding
narrow paper-claim term obtained from the existential it produces.

The module \path{Framework/StructuralLimitsEReal.lean} confirms the
mathematical correctness of the boundary values in extended-real
arithmetic using $\mathrm{ENNReal.log}$ (which sends $0$ to $\bot$
rather than the junk value $0$). The two modules together close the
gap: \texttt{GFullFix} is the Lean-level witness-family seal, and
\texttt{StructuralLimitsEReal} is the independent confirmation that
the value those witnesses produce is the mathematically correct one.

\subsection{Path~C$'$: full-real-domain trig}\label{sec:trig}

The three Path~C$'$ widenings deliver
\begin{itemize}[topsep=2pt, itemsep=1pt]
  \item \texttt{paper\_claim\_sin\_full} via
    $\sin x = \cos(\pi/2 - x)$;
  \item \texttt{paper\_claim\_arctan\_full} via
    $\arctan x = \arcsin(x / \sqrt{1 + x^2})$;
  \item \texttt{paper\_claim\_tan\_full} via the
    $\tan x = \tan(x - k\pi)$ period reduction.
\end{itemize}
\noindent
The shared infrastructure is the substitution operator
$\mathrm{subst}_0 : \EMLTC \to \EMLTC \to \EMLTC$ that replaces every
$\mathbf{x}_0$ in a host term with a given substitute term, together
with the foundational lemma $\mathrm{ADDsafe}_\C^{\,\mathrm{ofReal},\,
\mathrm{ofReal}}$ which discharges the 11-field
$\mathrm{ADDsafe}_\C$ precondition for $\mathrm{mkAdd}_\C$ whenever
both arguments are real-valued. This lets period shifts via repeated
$\mathrm{mkAdd}_\C$ of fixed real period constants stay entirely in
the real fragment, where the $\arg = \pi$ branch-cut trap never
appears.

\subsection{Closed numeric constants and the imaginary unit}

Five reusable \texttt{EMLRealization}$_\C$ packages
($\mathrm{realize}_\C^{\,\{0, 2, -\imu, \imu, \pi\}}$) live in
\path{Framework/Complex/Closures/Constants.lean}. Each consists of a
literal $\EMLTC$ tree plus a partial-evaluation proof that the tree
returns the constant when given any environment. These are the
building blocks for the trig witnesses, the Path~C$'$ period shifts,
and the §G \texttt{mkZero} term.

For the hand-tuned closed constants our K-counts match the paper's
hand-tuned Table~4 figures to the unit: $K(0) = 7$, $K(2) = 19$,
$K(-\imu) = 127$, $K(\imu) = 407$, $K(\pi) = 233$.

\subsection{Sheffer cousins: EDL and $-$EML}\label{sec:sheffer}

The paper presents EML, EDL, and $-$EML as a ``family'' of Sheffer
operators (paper §3.1, equation block \texttt{label\{Sheffers\}}) but
proves completeness only for EML; the cousins are confirmed
empirically via the Mathematica/Rust \texttt{VerifyBaseSet} procedure.

The artefact's \path{Framework/Sheffer.lean} scaffolds the two
paper-named cousins as inductive grammars with partial-evaluation
semantics:

\begin{itemize}[topsep=2pt, itemsep=2pt]
  \item \texttt{EDLTerm} with the binary
    $\edl(x, y) = \exp(x) / \log(y)$ paired with the constant $e$;
  \item \texttt{NegEMLTerm} with the binary
    $-\eml(y, x) = \log(x) - \exp(y)$ paired with the constant
    $-\infty$, which requires a parallel \texttt{NegEMLTermE} grammar
    over $\mathrm{EReal}$ to seal the $-\infty$ constant itself.
\end{itemize}

Eight of the $36$ EDL primitives and five of the $36$ $-$EML
primitives are sealed in this module (Plans D and E in
\path{OPEN_QUESTIONS.md}). The non-trivial EDL witnesses include
Aristotle's three-step composition for $\log x$ and the
$\edl(\mathrm{D8}(x), \mathrm{D4}(y))$ construction for division.

\subsection{The conditional ceiling scaffold}\label{sec:ceiling}

\path{Framework/EDLClosedVal.lean} formalises a closure result for
the closed-EDL value set:

\begin{lstlisting}[language=Lean]
inductive EDLClosedVal : Real -> Prop where
  | one     : EDLClosedVal 1
  | e_const : EDLClosedVal (Real.exp 1)
  | edl     : {a b : Real} ->
              EDLClosedVal a -> EDLClosedVal b ->
              Real.log b /= 0 ->
              EDLClosedVal (Real.exp a / Real.log b)

theorem edl_closed_eval_in_closedVal :
    forall {t : EDLTerm}, t.IsClosed ->
    forall (env : Nat -> Real) {v : Real},
      t.eval? env = some v -> EDLClosedVal v
\end{lstlisting}

\noindent
The closure theorem is proved unconditionally by induction on the
$\EDLT$ structure. To turn it into per-primitive non-membership
statements (no closed EDL term evaluates to $-1$, $2$, or $\frac{1}{2}$),
we need three transcendence-style facts whose proof is currently
out of Mathlib's reach. We package them as a Lean typeclass:

\begin{lstlisting}[language=Lean]
class EDLTranscendenceBarrier : Prop where
  neg_one_not_closed : ~ EDLClosedVal (-1)
  two_not_closed     : ~ EDLClosedVal 2
  half_not_closed    : ~ EDLClosedVal ((1 : Real) / 2)
\end{lstlisting}

\noindent
Three obstruction corollaries (e.g. \texttt{no\_closed\_edl\_neg\_one})
take the typeclass as a hypothesis. We \emph{intentionally provide no
instance}, so the corollaries are scaffolded but not closed. A future
Mathlib Schanuel-style transcendence result would discharge the
typeclass and close them; until then, any user invoking these
corollaries is signing for the conjecture explicitly. This is the
``conditional ceiling scaffold'' framing used throughout the
artefact's public documentation.

%==========================================================================
\section{The frontier sprint}\label{sec:frontier}

Four research-grade directions were identified by an independent
GPT~Pro architectural review (separate context, no shared scratchpad)
in early May 2026. A focused two-day sprint on 2026-05-10/11 produced
substantive Lean content for all four, adding $39$ public theorems.

\subsection{Direction 1: variable-transplant identity depths
(SI~§1.5~\#5)}

\noindent\textbf{Module:} \path{Framework/TransplantDepths.lean}
($9$ theorems).

The paper's identity term has depth four; the Supplementary Information
asks (open question \#5) whether other depths admit identity terms.
\texttt{TransplantDepths.lean} proves:

\begin{itemize}[topsep=2pt, itemsep=2pt]
  \item \textbf{Affirmative family:} for every $k \in \N$, an explicit
    $\EMLT$ term constructor \texttt{identityTermAtDepth\,k} of depth
    $4k$ that evaluates to $\mathrm{env}\,0$ on every real environment.
  \item \textbf{Negative closures for $d = 1$ and $d = 2$:} no
    $\EMLT$ of either depth is identity on every environment.
  \item \textbf{$d = 3$ statement:} the proposition
    \texttt{NoIdentityAtDepthThree\_conjecture\,:\,Prop}, recording
    the next open case. Aristotle proved it in a simplified
    single-atom grammar
    (\path{process_archive/chunks/090_no_identity_depth_3/result.lean});
    the canonical-grammar port (our two-atom $\{1, \mathrm{var}\}$
    grammar) is the natural follow-up.
\end{itemize}

\subsection{Direction 2: §G boundary points}

Covered in \S\ref{sec:gfix}.

\subsection{Direction 3: Plan~D conditional ceiling}

Covered in \S\ref{sec:ceiling}.

\subsection{Direction 4: paper~§5 universal minimality,
polynomial-class first cut}

\noindent\textbf{Module:} \path{Framework/PolynomialBinary.lean}
($2$ theorems).

The paper~§5 universal-minimality question asks whether \emph{any}
single-binary-operator scaffold $\{c, B\}$ can match EML's coverage.
The general question is open; the polynomial-class first cut is
fully proved.

\begin{definition}[Polynomial binary]
A binary operation $B \colon \R \to \R \to \R$ is \emph{polynomial}
when there exists a bivariate polynomial witness
$P \in \mathrm{MvPolynomial}(\mathrm{Fin}\,2,\,\R)$ such that
$B(x, y) = P(x, y)$ for all $x, y \in \R$.
\end{definition}

\begin{lemma}[Composition lemma]
If $B$ is polynomial and $t$ is any free term over $B$ with constants
and a single variable, then there exists a univariate polynomial
$Q \in \R[X]$ such that $t.\mathrm{eval}\,B\,(\mathrm{const}\,x) =
Q(x)$ for every $x \in \R$.
\end{lemma}

\begin{theorem}[Polynomial-binary impossibility]
\label{thm:polybinary}
No polynomial binary operation, freely composed with constants and a
single variable, can equal $\mathrm{Real.exp}$ on every input.
\end{theorem}

\begin{proof}[Proof sketch]
By the composition lemma, the candidate witness reduces to a
univariate polynomial $Q$ with $Q(x) = \exp(x)$ on every $x$.
Then $Q(x) / \exp(x) = 1$ identically. But
\texttt{Polynomial.tendsto\_div\_exp\_atTop} says this ratio tends
to $0$ at infinity, contradicting the constant-$1$ identification.
\end{proof}

This is a single class of rule-outs. Rational, semialgebraic, Pfaffian,
and real-analytic classes each need their own impossibility argument,
which remain paper-open.

\subsection{Direct-macro alternative witnesses: an honest finding}

\noindent\textbf{Module:} \path{Framework/AlternativeWitnesses.lean}
($9$ alternatives + $9$ K-counts).

For the seven non-trivial binary primitives we built a parallel
``direct-macro'' set of witnesses bypassing the structural compiler.
The expectation was K-count reduction; the actual finding is that the
direct-macro witnesses have \emph{identical} K-counts to the
compile-output ones, because the structural compiler already uses
these macros internally. The high K-counts reflect the genuine cost
of full-natural-domain coverage, not a compiler inefficiency.

The module was named \texttt{CompactWitnesses.lean} during development
and renamed to \texttt{AlternativeWitnesses.lean} after the finding to
avoid misleading future readers.

%==========================================================================
\section{Limitations and architectural findings}\label{sec:limits}

\subsection{Plan~B: custom complex-log branch is architecturally
infeasible}

The paper at line~333 sketches an alternative route to full real-domain
trig coverage by ``manually correcting the $i$-sign'' inside the
compiler. The natural Lean translation is to introduce a custom
branch-of-log primitive in $\EMLTC$ and route certain trig witnesses
through it.

This route foundered on a structural fact: $\EMLTC$'s partial
evaluation (\texttt{EMLTerm}$_\C$.\texttt{eval?} in
\path{Framework/Complex/Term.lean}) hard-codes Mathlib's principal
\texttt{Complex.log}. There is no place in the grammar definition
to substitute a different branch convention. Adding one would require
either (a) extending the grammar with a parameterised
$\mathrm{mkLog}_{\C, \theta}$ constructor (and propagating the
parameter through every macro), or (b) reinterpreting the existing
constructor through a custom evaluator (which breaks every existing
\texttt{eval?} lemma).

We concluded the route was not worth the architectural cost, and
adopted Path~C$'$ (range-reduction by substitution) instead. Path~C$'$
recovers the same full-real-domain coverage with no grammar change,
at the cost of moving from single-witness statements to
witness-family statements. The detailed finding is recorded in
\path{OPEN_QUESTIONS.md} §B.0.

\subsection{Conjectural items intentionally left open}

\begin{itemize}[leftmargin=*, topsep=2pt, itemsep=2pt]
  \item \textbf{EDL transcendence barrier.} The three obstruction
    corollaries for $-1$, $2$, $\frac{1}{2}$ are conditional on
    \texttt{EDLTranscendenceBarrier} (\S\ref{sec:ceiling}). No
    instance is provided.
  \item \textbf{Universal minimality beyond polynomial.}
    Theorem~\ref{thm:polybinary} rules out polynomial binaries only.
    The rational, semialgebraic, Pfaffian, and real-analytic classes
    each remain open.
  \item \textbf{Variable-transplant depth 3.} The
    canonical-grammar non-existence at $d = 3$ is recorded as a
    \texttt{Prop} (\texttt{NoIdentityAtDepthThree\_conjecture}).
    Aristotle has the proof in a simplified single-atom grammar.
  \item \textbf{SI §1.5 questions 1, 2, 3, 4, 6, 7.} The paper's own
    open-questions list. Out of artefact scope.
  \item \textbf{§4.3 gradient-based symbolic regression.}
    Out of artefact scope (Mathlib has no gradient-flow / projection
    / floating-point $\leftrightarrow$ symbolic infrastructure).
\end{itemize}

%==========================================================================
\section{Verification evidence}\label{sec:verification}

\subsection{Build}

\texttt{lake build EML} finishes successfully at $8062$ jobs, with
warnings confined to two cosmetic linters
(\texttt{linter.unusedSimpArgs}, \texttt{linter.unusedVariables}). A
verbatim build-tail transcript is in
\path{lambda_lab/proofs/eml/2603_21852/VERIFICATION_EVIDENCE.md}.

\subsection{Axiom audit}

The audit script \path{lean_workspace/EML/AxiomCheck.lean} runs
\texttt{\#print axioms} on every headline theorem. The verbatim
output (also in \path{VERIFICATION_EVIDENCE.md}) shows that every
public theorem depends only on Mathlib's three standard axioms:
\texttt{propext}, \texttt{Classical.choice}, \texttt{Quot.sound}.
The \texttt{K\_count\_*} theorems go further: they are pure
\texttt{rfl}-decidable and depend on \emph{no} axioms at all.

\noindent\textbf{What this rules out.}
\begin{itemize}[topsep=2pt, itemsep=2pt]
  \item No \texttt{sorry} hides inside the chain (the axiom set is
    the closed Mathlib-standard one).
  \item No project-side \texttt{axiom} declaration was introduced.
  \item The \texttt{EDLTranscendenceBarrier} typeclass, the one
    place where the artefact \emph{consciously} uses an unproved
    hypothesis, does not appear in the axiom set of the
    unconditional closure theorem
    \texttt{edl\_closed\_eval\_in\_closedVal}. The three conditional
    corollaries take it as a parameter and would print it in their
    axiom list if asked; by design no instance is provided.
\end{itemize}

\subsection{Repository hygiene}

\texttt{grep -rn 'sorry|admit' lean\_workspace/EML/ --include=*.lean}
returns nothing matching an actual tactic site. The four false-positive
matches are all comment text describing the word ``sorry'' (e.g.
``a permanent \texttt{sorry}'' from the phase-1 report). The CI hook
\path{scripts/check_no_sorries.sh} enforces this on every commit.

\subsection{Independent re-verification}

PCSS Eagle HPC re-verifies the artefact on a different machine and a
different Lean toolchain installation. Most recent run: job
$7\,041\,555$ on 2026-05-07, $88$ files, $0$ failures, $42$\,s wall
time. The submission scripts are in \path{eagle_scripts/}.

%==========================================================================
\section{Reproducing the verification}\label{sec:repro}

\subsection{Local}

\begin{lstlisting}[language=bash, basicstyle=\ttfamily\small]
git clone https://github.com/nasqret/eml-formalization
cd eml-formalization/lambda_lab/proofs/eml/2603_21852/lean_workspace
lake build EML                       # ~30-60 min cold cache
lake env lean EML/AxiomCheck.lean    # axiom audit
\end{lstlisting}

\noindent
The first build pulls $\sim 6$\,GB of Mathlib oleans and finishes in
$30$--$60$ minutes; subsequent builds are incremental.

\subsection{On PCSS Eagle (HPC)}

\begin{lstlisting}[language=bash, basicstyle=\ttfamily\small]
# From a node where ~/pl0414-02/scratch/elan_dl/
# has the Lean toolchain tarball
cd /mnt/storage_5/scratch/pl0414-02
sbatch verify_all.sbatch    # parallel verify across 24 cores
\end{lstlisting}

\noindent
The SLURM scripts are in \path{eagle_scripts/} and handle (a) Lake
olean cache rebuild and (b) parallel re-verify.

%==========================================================================
\section{Where the formalisation diverges from the paper}\label{sec:diverge}

Three architectural moves separate our artefact from a literal
transcription of the paper's prose.

\begin{enumerate}[leftmargin=*, topsep=2pt, itemsep=4pt]
  \item \textbf{Partial evaluation as junk-value avoidance.} The paper
    treats $\log 0 = -\infty$ as routine; we make every nested
    $\eml$-evaluation return \texttt{none} outside the natural domain
    of $\log$. The bridge theorems are stated as ``if
    $\F$-eval returns \texttt{some\,v}, then there exists an $\EMLT$
    term that also returns \texttt{some\,v}''; we never claim
    equality at a boundary point.
  \item \textbf{Three §G boundary points as witness families.} The
    paper's witness for $\sqrt{0}$, $\mathrm{arcosh}\,1$, and
    $\mathrm{hypot}(0, 0)$ would implicitly use the $\log 0 = -\infty$
    convention; Mathlib's $\log 0 = 0$ blocks this. We use the
    quantifier flip described in \S\ref{sec:wf}. The paper itself
    (line~342 of \texttt{EML.tex}) explicitly notes this Lean-specific
    issue.
  \item \textbf{Path~C$'$ instead of the paper's ``manual $i$-sign
    correction''.} The paper's prose at line~333 sketches a custom
    branch of $\log_\C$. We found this architecturally infeasible in
    Lean (\S\ref{sec:limits}) and used range-reduction by substitution
    instead. Both approaches deliver full-real-domain trig coverage;
    the paper's is implicit in compiler code that chooses different
    witness shapes per region of $x$, ours is explicit in the
    quantifier flip.
\end{enumerate}

%==========================================================================
\section{Conclusion and acknowledgements}\label{sec:conclusion}

The Lean 4 + Mathlib v4.28 artefact described in this report seals all
$36$ primitives of the paper's reference calculator via literal
$\EMLT$ or $\EMLTC$ witness terms on a non-empty open subdomain of
their natural mathematical domain. The three §G boundary points the
paper itself flags are additionally sealed by witness-family theorems
in which the witness depends on the environment. After the
post-submission frontier sprint, the public API exposes $100$
machine-checked theorems covering the original paper claims,
full-domain trig, §G boundary fixes, variable-transplant identity
depths, and a polynomial-class first cut of paper~§5 universal
minimality. The build is sorry-free at $8062$ jobs and every public
theorem depends only on Mathlib's three standard axioms.

The architectural shifts that made this possible are recorded in
detail in \path{AUTHOR_SUMMARY.md} and
\path{notes/proof_structure.pdf}; the live status of every open item
is in \path{OPEN_QUESTIONS.md}; the fresh re-verification evidence is
in \path{VERIFICATION_EVIDENCE.md}.

\subsection*{Acknowledgements}

\begin{itemize}[leftmargin=*, topsep=2pt, itemsep=1pt]
  \item \textbf{Andrzej Odrzywo\l{}ek} (Jagiellonian University) ---
    the source paper. Thanks for both the discovery of the EML
    operator and for the careful description of the §G boundary issue
    (paper line 342) which spared us a great deal of confusion when
    we first hit it in Lean.
  \item \textbf{Aristotle} (Harmonic AI) --- proof search for many
    individual chunks during the phase-1 chunked decomposition and
    the frontier sprint. The full chunk archive is in
    \path{process_archive/chunks/}.
  \item \textbf{GPT~Pro} (OpenAI) --- three independent-context
    architectural reviews. Recommended the structural-compiler
    architecture, the Cayley-quotient route for $\tan$, the public
    closed-constants packaging, Path~C$'$, the witness-family
    quantifier flip, the four frontier-sprint directions, and the
    pre-announcement punch list. Transcripts in
    \path{process_archive/gpt_pro_bundle/}.
  \item \textbf{Claude Code} (Anthropic) --- orchestration,
    scaffolding, integration of Aristotle returns, post-submission
    trig widenings, and frontier-sprint module assembly.
  \item \textbf{Codex} (OpenAI) --- paraphrase and informalisation
    of mathematical statements during paper-to-grammar translation.
  \item \textbf{Mathematica} --- enumeration and witness candidate
    search during phase-1 decomposition.
  \item The \textbf{Mathlib community} --- the underlying Lean
    library.
\end{itemize}

\noindent\textbf{Phase-1 process history.} An earlier
$60$-page report dated 27 April 2026, describing the chunked
decomposition approach (since superseded by the structural compiler
and the frontier sprint), is preserved at
\path{process_archive/REPORT_2026_04_27.pdf} for the audit trail.

\end{document}
